\documentclass{esannV2}
\usepackage[utf8]{inputenc}
\usepackage{graphicx}
\usepackage{amssymb,amsmath,array}
\usepackage{booktabs}

% Margins are fixed by esannV2.cls: 5.2 cm top/bottom, 4.4 cm left/right on A4.
% Do not change them. Text block 12.2 cm; the figure was saved at exactly
% that width.
\voffset 0 cm \hoffset 0 cm \addtolength{\textwidth}{0cm}
\addtolength{\textheight}{0cm}\addtolength{\leftmargin}{0cm}

\begin{document}

\title{Participation ratio flags some degenerate representations,\\but does not reliably rank healthy encoders}

% TODO: confirm the full name and the exact unit name before submission.
\author{Guilherme Lima
\vspace{.3cm}\\
Universidade Federal de Pelotas - Instituto de Informática \\
Pelotas, RS - Brazil
}

\maketitle

\begin{abstract}
The geometry of a frozen representation is often offered as a label-free way to
pick an encoder, on the premise that a measure correlated with downstream
accuracy can rank candidates. We test that premise on seventeen encoders across
four domains, in a cross design: nine architectures under one supervised
ImageNet regime, nine pretrainings over the same ViT-B/16. The participation
ratio does correlate with linear-probe accuracy ($\rho = 0.41$, $p = 0.007$),
but the correlation lives in the collapsed tail --- $0.27$ without MAE, $0.21$
without the two CLIP as well, $0.11$ across the nine architectures --- and
holding width constant leaves $+0.39$ ($p = 0.010$). A held-out checkpoint,
predicted before we measured it, shows the premise failing where it matters:
BEiT-3 ranks among the worst of the eighteen on the probe, yet the ratio sees
the problem in only one domain of four. The measure flags some degenerate representations, but we find no reliable
evidence that it ranks healthy ones, and it suits diagnosis better than
selection.
\end{abstract}

\section{Introduction}

Choosing a frozen encoder for a small domain is a measurement problem. The
usual answer, comparing downstream accuracy, is mostly noise at that sample
size: resampling the partition moves the score by more than most architectural
differences \cite{Bouthillier}, and cross-validation error bars at $n \approx
100$ reach $\pm 10$ points \cite{Varoquaux}. In a recent hematology study a
source-fitted linear probe predicted only loosely which frozen encoder would
survive a change of equipment, with median $\rho = 0.45$ between the two
orderings \cite{Hemato}. Geometry of the frozen representation promises a way
out: it is label-free and costs one forward pass per encoder, so it would let a
practitioner choose without running the evaluation that is too uncertain to
trust. The promise rests on a step that is rarely tested, namely that a measure
\emph{correlated} with accuracy can \emph{rank} candidates. We test that step
for the participation ratio.

Testing it needs a design that separates architecture from pretraining. When
both vary together, as in most encoder comparisons, a geometric difference
between two models admits at least two explanations and the data cannot choose
between them: an apparent architecture effect may be the pretraining recipe
each family ships with. Our cross design therefore holds the pretraining
\emph{regime} fixed, supervised ImageNet classification, across nine
architectures whose recipes still differ, and varies nine pretrainings over one
ViT-B/16 architecture and parameterisation. The two arms meet at a shared
checkpoint, and the seventeen encoders cross four domains, giving 68
encoder--domain pairs.

Two contributions. We locate where the association between participation ratio
and probe accuracy actually lives --- in the collapsed tail, not in a fine
ordering --- and show width does not explain it. And we test the resulting
claim out of sample, on held-out checkpoints with the prediction registered
beforehand, which bounds it from both sides.

\section{Measures and protocol}

\subsection{Encoders and domains}

Architecture arm, all pretrained by supervised ImageNet classification but
under different recipes and from different providers: ResNet-18 and -50
\cite{ResNet}, ConvNeXt-T
\cite{ConvNeXt}, EfficientNet-B3 \cite{EffNet}, Swin-T \cite{Swin}, MaxViT-T
and -S \cite{MaxViT}, ViT-S/16 and ViT-B/16 \cite{ViT}. Pretraining arm, the
same ViT-B/16 in nine versions: four supervised with different corpora and
recipes (IN-21k$\to$1k, IN-21k, IN-1k, DeiT3 \cite{DeiT3}), two
self-supervised (DINO \cite{DINO}, MAE \cite{MAE}) and three with language
(CLIP on two corpora \cite{CLIP}, SigLIP \cite{SigLIP}).

Four domains, greyscale and $512\times512$ so colour is not a confound --- the
single channel is replicated across the three inputs and each checkpoint keeps
\emph{its own} pretraining normalisation: \textbf{mammography} (INbreast
\cite{INbreast}), \textbf{DTD} \cite{DTD}, \textbf{Pet} \cite{Pet} and
\textbf{SUN397} \cite{SUN397}, with 2, 47, 37 and 22 classes; 387 images each
except SUN397, 384 after dropping classes with fewer than five examples. Images
were drawn round-robin across classes under a fixed seed, so the sample is
near-balanced where the dataset allows (8--9 per class on DTD, 5--16 on Pet,
3--24 on SUN397); the distribution matters because the participation ratio
comes from the global covariance, which mixes within- and between-class
variation. The SUN397 crop is one shard of the public mirror, so the accuracies
are comparable between our encoders but should not be read against the transfer
literature.

No encoder is fitted; the probe below is the only trained component. Geometry
uses the first 200 images per domain, the probe the full sample. In the ViTs at
512\,px the positional embeddings are interpolated to the 32$\times$32 grid
(\texttt{timm}'s \textsc{dynamic\_img\_size}); the MaxViTs are native
512\,px checkpoints. Checkpoint identifiers, seeds, extraction and permutation
code, the registered prediction of Section 3.2 and the full per-checkpoint
table are in the
repository.\footnote{\textsc{repository url --- fill in before submission}}

\subsection{The measures}

\paragraph{Extraction.} Identical across all seventeen: the last feature map,
then \textbf{global average pooling}, with ViT prefix tokens \emph{discarded}
and spatial tokens regridded first --- the \textsc{cls} token is never used, so
convolutional networks and transformers do not differ in extraction strategy.

\paragraph{Participation ratio.} $(\sum_i \lambda_i)^2/\sum_i \lambda_i^2$ over
the covariance eigenvalues of that representation: how many directions it
actually uses. It is \emph{not} the entropic effective rank \cite{EffRank}, and
we prefer it because it is less swayed by the smallest eigenvalues.

\paragraph{Linear probe.} Standardised logistic regression, five outer folds,
stratified by class and \textbf{grouped by patient} on mammography (103
patients), with $C$ searched over three inner folds per encoder, domain and
fold; fixing it would measure width instead. Scores come from the held-out fold
(AUC on mammography, accuracy elsewhere) and need not match across domains,
because every correlation below uses scores standardised inside each domain. Three independent partitions (seeds 0, 1, 2)
differ by 0.002--0.006, far below the interval widths.

\subsection{Uncertainty}

Intervals come from resampling \emph{images}, and \emph{patients} on
mammography; for the probe the resamples of the three partitions are pooled.
The participation ratio is the exception. Drawing with replacement leaves
fewer distinct images and moves the spectrum for a reason that has nothing to
do with the variability we want, so we subsample \emph{without} replacement at
$m = n/2 = 100$, drawing 400 subsamples and reporting point and interval at the
same $m$. That interval is the 2.5th and 97.5th percentile of the 400 values,
unadjusted, and so estimates how much the ratio moves across 100-image subsets
of this sample --- not uncertainty about a population value, which would need a
model of how the sample arose. The covariance of $m$ images has rank at most
$m-1$, capping the estimate at 99.

Correlation points are \emph{not} independent: each encoder appears in all four
domains. Inference permutes encoder labels, \textbf{the same permutation across
all four domains} (20\,000 draws); shuffling each domain on its own would make
the null too random and the test too easy to pass. Intervals come from a
leave-one-encoder-out \emph{jackknife}, taken after Fisher's z-transform.

\section{Results}

\subsection{The ordering is stable, the magnitudes are not}

How much of a measure's spread comes from the dataset rather than the encoder
depends heavily on the measure: 43\,\% for the participation ratio, 12\,\% for
how far a central output position looks back into the image (a gradient-based
radius in input pixels \cite{Luo}), and 0.5\,\% for how much each transformer
block changes what it receives, $\|F(x)\|/\|x\|$, on the nine ViT-B/16 where
blocks line up --- the last two are defined in full in the repository. Two
orders of magnitude separate the extremes, and the order holds on sets of 6, 12
and 17 encoders, which are simply the measurement batches as they accumulated,
so no selection is involved. What does not hold is the size of the numbers:
that first share was 66\,\%, 31\,\% and 43\,\% on the three sets. A share
like this describes the set it was measured on and little else.

Sample size behaves the same way, and the ceiling of 99 is not remote: the
largest ratio we measure is 55.4, the only one of 68 pairs above half, with a
median of 20.1. Magnitudes depend on $m$ --- the median of PR(50)/PR(200) is
0.74 --- but the ordering does not. Each size being the mean of 40 subsamples,
Spearman between the ratio at 50 and at 200 images is 0.990--1.000 in all four
domains (one comparison per domain, over the means), and
Table~\ref{tab:sens} is indistinguishable on half the images. The collapsed
checkpoint is the steadiest: 1.81 at both 50 and 200 images.

\subsection{Where the association comes from}

Scoring every encoder on both and standardising inside each domain gives
$\rho = +0.41$ over the seventeen ($p_{\text{perm}} = 0.007$; 95\,\%
CI $[+0.07; +0.67]$). The association weakens sharply once extreme checkpoints
leave: Table~\ref{tab:sens} shows $\rho$ falling to
$+0.27$ without MAE --- the only collapsed checkpoint, 1.8 to 4.2 directions
--- and to $+0.21$ without the two CLIP as well, with the interval containing
zero. Across the nine architectures, none collapsed, $\rho = +0.11$. The
intervals still leave room for a moderate association, reaching $+0.51$ at the
upper end, and we ran no equivalence test against a smallest effect of
interest. We can say that no reliable evidence of ordering ability
turned up, but not that the ability is absent.

Width in channels, the obvious confound, is associated with the probe on its
own ($+0.17$); controlling for it by partial Spearman leaves the ratio at
$+0.39$ ($p = 0.010$). The partial is the correlation of the residuals of the
ratio and probe ranks after both are regressed on width ranks. Its $p$ needs
care, because the hypothesis to rule out is that the ratio adds nothing once
width is known, and a null for it has to keep the width--probe link alive
rather than break it. We therefore use Freedman--Lane: regress the probe on
width, shuffle only the residuals (at encoder level, same permutation across
domains) and add the untouched fit back. In the
equivalent regression, standardised within domain, the width coefficient is
$+0.02$ against $+0.57$ for the ratio ($p = 0.003$, same null). The field
radius, measured throughout alongside the participation ratio, shows no
association with the probe ($+0.20$, $p = 0.26$).

\begin{table}[t]
\centering\small
\setlength{\tabcolsep}{4pt}
\begin{tabular}{@{}lrrrc@{}}
\toprule
set & $k$ & $\rho$ & $p$ & 95\,\% CI \\
\midrule
all 17 & 17 & $+0.41$ & 0.007 & $[+0.07; +0.67]$ \\
without MAE & 16 & $+0.27$ & 0.075 & $[-0.01; +0.51]$ \\
without MAE and both CLIP & 14 & $+0.21$ & 0.223 & $[-0.10; +0.48]$ \\
\midrule
9 pretrainings & 9 & $+0.41$ & 0.088 & $[-0.18; +0.78]$ \\
9 pretrainings, without MAE & 8 & $+0.26$ & 0.273 & $[-0.16; +0.60]$ \\
9 architectures & 9 & $+0.11$ & 0.662 & $[-0.34; +0.51]$ \\
\bottomrule
\end{tabular}
\caption{Sensitivity of $\rho$ (participation ratio against probe) to dropping
the checkpoints of extreme dimensionality. $p$ by encoder permutation, interval
by leave-one-encoder-out jackknife on the Fisher scale; the interval contains
zero exactly where $p > 0.05$.}
\label{tab:sens}
\end{table}

We call a checkpoint \emph{collapsed} only after looking: MAE's ratio is
1.8--4.2, three to ten times below the next encoder in any domain, and it is
the worst probe of the seventeen in all four (0.533 AUC; 0.391, 0.322, 0.627
accuracy). That is circular and we should say so --- low ratio is exactly what
we mean by collapsed, so inside this set ``the ratio flags collapse'' cannot be
proved wrong.

We tested outside it, on two \emph{held-out} encoders measured by the same
protocol, the prediction registered in a repository file predating the first
recorded result. The flag rule is fixed in advance and applied per domain: an
encoder is \emph{flagged} where its ratio falls below the lowest non-collapsed
reference value in that domain (4.93 mammography, 9.78 DTD, 11.35 Pet, 10.43
SUN397), so collapse here describes the geometry within a domain rather than
the checkpoint as a whole. An untrained ViT-B/16 is flagged in all four, 3.9 to
6.8 times below that reference, and probes at chance. BEiT-3 base \cite{BEiT3}, a
masked-modelling checkpoint without classification fine-tuning, is the
informative case: its probe is second worst of eighteen on DTD, Pet and SUN397
and among the five worst on mammography, yet the ratio flags it on mammography
alone (2.2 against 4.9) and sits inside the distribution of the seventeen
elsewhere (17.7--22.8). The prediction failed in three domains of four, in the
direction that matters: there are poor representations that this geometry does
not see. Ranking low on both measures in every domain, that checkpoint
also carries the correlation in MAE's place ($\rho = +0.36$, $p = 0.018$,
against $+0.27$ with neither), though that is a secondary, unregistered
analysis.
Two encoders cannot pin down how often the measure is right, so what the test
gives us is a counterexample rather than a characterised detector.

Nine tests were run here; only the first and the held-out test were
pre-specified, and no multiplicity correction is applied. One exploratory
association, between field radius and probe in the pretraining arm
($\rho = +0.56$, $p = 0.009$), carries no confirmatory value.

\section{Discussion and limitations}

A representation can have healthy global dimensionality and still be poor for
the task. That is what the held-out BEiT-3 shows, and it is the substantive
lesson: geometry of this kind tells you when a representation has gone
degenerate, but not how good one is. The participation ratio does flag
degenerate representations in the domain where it is measured, by a margin of 2
to 7 times on the three collapsed checkpoints we found, though with no
guarantee of coverage, and we find no reliable evidence that it ranks healthy
encoders. That does not require the ability to be zero, only that seventeen
encoders could not demonstrate it. The field radius serves neither purpose, and
one point generalises beyond both measures: a number like ``43\,\% of the
spread comes from the dataset'' should be reported with the models it was
computed on.

Limitations. The probe is linear and frozen, and MAE is the reminder that low
dimensionality says the information is not linearly accessible, and says
nothing about whether it is there at all. In the pretraining arm the permutation test only calls a result
significant above $|\rho| = 0.46$, which is the null's 95th percentile rather
than a power calculation, so the $0.41$ there falls short. The held-out test covers one
collapsed checkpoint at a time, never a \emph{set} of them. A fourth measure,
patch-shuffling sensitivity, was discarded as a grid-phase artefact; it and a
greyscale control are in the repository. Finally, every measure here looks at
the representation as a whole and cannot see how its parts combine.

\begin{footnotesize}
\begin{thebibliography}{99}

\bibitem{Bouthillier} X. Bouthillier et al., Accounting for variance in machine learning benchmarks. \emph{MLSys}, 2021.

\bibitem{Varoquaux} G. Varoquaux, Cross-validation failure: small sample sizes lead to large error bars. \emph{NeuroImage}, 180:68-77, 2018.

\bibitem{Hemato} J. K. Sharma and P. Tapadiya, Can you trust frozen hematology foundation models under acquisition shift? arXiv:2608.25148, 2026.

\bibitem{INbreast} I. C. Moreira et al., INbreast: toward a full-field digital mammographic database. \emph{Academic Radiology}, 19(2):236-248, 2012.

\bibitem{DTD} M. Cimpoi, S. Maji, I. Kokkinos, S. Mohamed and A. Vedaldi, Describing textures in the wild. \emph{CVPR}, 3606-3613, 2014.

\bibitem{Pet} O. M. Parkhi, A. Vedaldi, A. Zisserman and C. V. Jawahar, Cats and dogs. \emph{CVPR}, 3498-3505, 2012.

\bibitem{SUN397} J. Xiao, J. Hays, K. A. Ehinger, A. Oliva and A. Torralba, SUN database: large-scale scene recognition from abbey to zoo. \emph{CVPR}, 3485-3492, 2010.

\bibitem{EffRank} O. Roy and M. Vetterli, The effective rank: a measure of effective dimensionality. \emph{EUSIPCO}, 606-610, 2007.

\bibitem{Luo} W. Luo, Y. Li, R. Urtasun and R. Zemel, Understanding the effective receptive field in deep convolutional neural networks. \emph{NIPS}, 4898-4906, 2016.

\bibitem{MaxViT} Z. Tu et al., MaxViT: multi-axis vision transformer. \emph{ECCV}, 459-479, 2022.

\bibitem{ViT} A. Dosovitskiy et al., An image is worth 16x16 words: transformers for image recognition at scale. \emph{ICLR}, 2021.

\bibitem{ResNet} K. He, X. Zhang, S. Ren and J. Sun, Deep residual learning for image recognition. \emph{CVPR}, 770-778, 2016.

\bibitem{Swin} Z. Liu et al., Swin transformer: hierarchical vision transformer using shifted windows. \emph{ICCV}, 10012-10022, 2021.

\bibitem{ConvNeXt} Z. Liu et al., A ConvNet for the 2020s. \emph{CVPR}, 11976-11986, 2022.

\bibitem{EffNet} M. Tan and Q. Le, EfficientNet: rethinking model scaling for convolutional neural networks. \emph{ICML}, 6105-6114, 2019.

\bibitem{DeiT3} H. Touvron, M. Cord and H. Jegou, DeiT III: revenge of the ViT. \emph{ECCV}, 516-533, 2022.

\bibitem{DINO} M. Caron et al., Emerging properties in self-supervised vision transformers. \emph{ICCV}, 9650-9660, 2021.

\bibitem{MAE} K. He et al., Masked autoencoders are scalable vision learners. \emph{CVPR}, 16000-16009, 2022.

\bibitem{BEiT3} W. Wang et al., Image as a foreign language: BEiT pretraining for all vision and vision-language tasks. \emph{CVPR}, 19175-19186, 2023.

\bibitem{CLIP} A. Radford et al., Learning transferable visual models from natural language supervision. \emph{ICML}, 8748-8763, 2021.

\bibitem{SigLIP} X. Zhai, B. Mustafa, A. Kolesnikov and L. Beyer, Sigmoid loss for language image pre-training. \emph{ICCV}, 11975-11986, 2023.

\end{thebibliography}
\end{footnotesize}

\end{document}
